\chapter{Friedreich's Ataxia}
\index{Friedreich's Ataxia}
\label{sec:Friedreich's Ataxia}

\section{Overview}
Friedreich's loss of coordination is the most common hereditary loss of coordination, caused by homozygous GAA trinucleotide repeat expansion in the FXN gene encoding frataxin, with a prevalence of 1 in 50,000. This neurological disorder affects the central or peripheral nervous system, leading to progressive or episodic symptoms. It significantly impacts quality of life and often requires multidisciplinary care. Neurological disorders affect the central or peripheral nervous system, leading to a wide range of symptoms affecting movement, sensation, cognition, and autonomic function. The nervous system's complexity means that even small disruptions can cause significant functional impairment. Many neurological conditions are chronic and progressive, requiring long-term management and rehabilitation. Advances in neuroimaging and molecular diagnostics have improved understanding and treatment of these conditions.
\section{Causes}
This condition happens because of frataxin deficiency, leading to mitochondrial iron accumulation, oxidative stress, and impaired energy metabolism, primarily affecting the dorsal root ganglia, posterior columns, corticospinal tracts, and cerebellum.    Protein aggregation and accumulation is a common mechanism in many neurodegenerative disorders.  Genetic mutations account for some neurological conditions, while others are sporadic or environmentally triggered.   
\section{Risk Factors}
\begin{itemize}[noitemsep]
  \item Genetic predisposition and family history
  \item Advancing age
  \item Lifestyle factors (diet, exercise, smoking, alcohol)
  \item Environmental exposures
  \item Underlying medical conditions
  \item Medication use
\end{itemize}
\section{Signs and Symptoms}
\subsection{Common symptoms}
\begin{itemize}[noitemsep]
  \item \item You may experience progressive gait and limb loss of coordination, slurred speech, areflexia, loss of vibratory and proprioceptive sensation, extensor plantar responses, scoliosis, pes cavus, hypertrophic cardiomyopathy, and increased risk of diabetes mellitus. Onset is typically in childhood or adolescence (mean age 10--15 years). Headache and facial pain Weakness or paralysis of muscles Numbness, tingling, or abnormal sensations Vision changes including blurred or double vision Difficulty with balance and coordination Cognitive changes (memory loss, confusion, difficulty concentrating) Speech and language difficulties Seizures or involuntary movements Dizziness and vertigo Fatigue and sleep disturbances
  \item Pain or discomfort in the affected area
  \end{itemize}
\subsection{Emergency warning signs}
\begin{itemize}[noitemsep]
  \item Severe or worsening symptoms of friedreich's ataxia require immediate medical evaluation
\end{itemize}
\section{Progression and Stages}
The condition may progress through different stages of severity over time. Early stages often involve mild or subtle symptoms that may be overlooked. Moderate stages involve more noticeable symptoms affecting daily function. Advanced stages may involve significant impairment and complications.
\section{Complications}
\begin{itemize}[noitemsep]
  \item Disease progression leading to worsening symptoms
  \item Secondary infections or injuries
  \item Side effects of treatment
  \item Reduced quality of life
  \item Emotional and psychological impact
\end{itemize}
\section{Diagnosis}
Doctors diagnose this based on clinical criteria and confirmed by FXN genetic testing. Neurological diagnosis combines clinical history and examination findings with neuroimaging (MRI, CT), electrophysiological studies (EEG, EMG, nerve conduction), and laboratory testing of blood and cerebrospinal fluid.

\begin{itemize}[noitemsep]
  \item Comprehensive medical history and physical examination
  \item Laboratory tests to assess organ function
  \item Imaging studies to visualize affected structures
  \item Specialized testing depending on the affected system
  \item Consultation with relevant specialists
\end{itemize}
\section{Prevention}
\begin{itemize}[noitemsep]
  \item Maintain a healthy lifestyle with balanced diet and exercise
  \item Avoid known risk factors
  \item Attend regular medical check-ups for early detection
  \item Follow recommended screening guidelines
\end{itemize}
\section{Treatment Options}
Treatment focuses on symptomatic treatment: physical therapy for balance and mobility, occupational therapy, speech therapy, and orthopedic Treatment for scoliosis. Cardiac monitoring is essential (echocardiography, ECG). Neurological treatment focuses on symptom management, disease modification when possible, and rehabilitation to maintain function. A multidisciplinary team approach is often beneficial. Pharmacologic therapy targeting specific neurotransmitter systems or disease mechanisms Physical, occupational, and speech therapy for functional maintenance Surgical interventions when appropriate (deep brain stimulation, tumor resection) Cognitive rehabilitation and memory support Pain management for neuropathic pain Assistive devices and mobility aids Lifestyle modifications including exercise, nutrition, and sleep hygiene Support services and caregiver education Regular neurologic monitoring and treatment adjustment Clinical trials for emerging therapies

\begin{itemize}[noitemsep]
  \item Medications to manage symptoms and address underlying mechanisms
  \item Lifestyle modifications and self-care measures
  \item Physical therapy and rehabilitation
  \item Surgical intervention when indicated
  \item Regular monitoring and follow-up care
\end{itemize}
\section{Living with It}
\begin{itemize}[noitemsep]
  \item Follow your treatment plan as prescribed
  \item Maintain regular follow-up with your healthcare provider
  \item Make necessary lifestyle adjustments
  \item Seek support from family, friends, or support groups
  \item Stay informed about your condition
\end{itemize}
\section{Outlook}
\begin{itemize}[noitemsep]
  \item Outlook is progressive
  \item Most patients lose ambulation 10--15 years after symptom onset, with mean survival to age 40--50, primarily due to cardiomyopathy.
\end{itemize}
